%% kapa: exam=2 round=28 year=2017 period=1 subject=appraisal_practice subject_ko=감정평가실무 seq=1 number=1 points=40
%% source: Q-Net 감정평가사 기출문제 (cst00309) artlSeq=5201712 file=raw/1008/5201712/2017년 감정평가사 2차시험 기출문제(감정평가 실무).pdf
%% conversion: kordoc 4.24.51 -> tools/gen2.py (auto); review=proofread-sonnet; flags=none
\begin{longquestion}{【문제 1】}
감정평가사 김○○은 W시 N구청장으로부터 도시계획시설(도로)사업과 \par
관련하여 토지의 보상평가를 의뢰받았다. 관련 법규 및 이론을 참작하고 제시된 \par
자료를 활용하여 다음 물음에 답하시오. (40점) \par
\medskip\noindent 물음 1) 미지급용지의 개념 및 평가기준을 기술하고, 대상토지의 감정평가액을 \par
구하시오. (15점) \par
\medskip\noindent 물음 2) 사실상 사도의 개념 및 평가기준을 기술하고, 대상토지의 감정평가액을 \par
구하시오. (15점) \par
\medskip\noindent 물음 3) 예정공도의 개념 및 평가기준에 대해 기술하고, 대상토지의 감정평가액을 \par
구하시오. (10점) \par
<공통 자료> \par
1. 사업의 개요 \par
1) 사업시행지 : W시 N구 M동 100-4번지 일원 \par
2) 사업의 종류 : 도시계획시설(도로)사업(소로2-60호선) 개설공사 \par
3) 사업시행자 : W시 N구청장 \par
4) 사업의 착수 예정일 및 준공예정일 : 인가일 ～ 2018.03.31 \par
2. 사업추진일정 \par
{\small\centering
\begin{tabular}{|c|c|}
\hline
\makecell{구분} & \makecell{일정} \\
\hline
\makecell{ㆍ도시계획시설(도로)결정일} & \makecell{2010.08.21} \\
\hline
\makecell{ㆍ도시계획시설(도로)사업 실시계획인가고시일} & \makecell{2016.12.15} \\
\hline
\makecell{ㆍ보상계획공고일} & \makecell{2017.03.31} \\
\hline
\makecell{ㆍ현장조사완료일} & \makecell{2017.06.01} \\
\hline
\end{tabular}\par}

※ 보상의뢰서상 가격시점 요구일 : 2017.07.01 \par
3. 대상토지의 개요(가격시점 현재) \par
{\small\centering
\begin{tabular}{|c|c|c|c|c|c|c|}
\hline
\makecell{기호} & \makecell{소재지} & \makecell{편입면적\newline{}(m\textsuperscript{2})} & \makecell{지목} & \makecell{현실\newline{}이용상황} & \makecell{용도지역} & \makecell{비고(소유자)} \\
\hline
\makecell{1} & \makecell{W시 N구 M동\newline{}100-4번지} & \makecell{381} & \makecell{전} & \makecell{도로} & \makecell{준주거지역} & \makecell{홍길동} \\
\hline
\end{tabular}\par}

4. 표준지 공시지가 자료 \par
{\small\centering\footnotesize\setlength{\tabcolsep}{3pt}
\begin{tabular}{|c|c|c|c|c|c|c|c|c|c|}
\hline
\multirow{2}{*}{\makecell{기호}} & \multirow{2}{*}{\makecell{소재지}} & \multirow{2}{*}{\makecell{면적\newline{}(m\textsuperscript{2})}} & \multirow{2}{*}{\makecell{지목}} & \multirow{2}{*}{\makecell{이용상황}} & \multirow{2}{*}{\makecell{용도지역}} & \multirow{2}{*}{\makecell{도로\newline{}교통}} & \multirow{2}{*}{\makecell{형상\newline{}지세}} & \multicolumn{2}{|c|}{\makecell{공시지가(원/m\textsuperscript{2})}} \\
\cline{9-10}
 &  &  &  &  &  &  &  & \makecell{2016.01.01} & \makecell{2017.01.01} \\
\hline
\makecell{A} & \makecell{W시 N구\newline{}M동\newline{}105번지} & \makecell{400} & \makecell{대} & \makecell{주거\newline{}나지} & \makecell{2종\newline{}일주} & \makecell{세로\newline{}(불)} & \makecell{가장형\newline{}평지} & \makecell{770,000} & \makecell{860,000} \\
\hline
\makecell{B} & \makecell{W시 N구\newline{}M동\newline{}103번지} & \makecell{420} & \makecell{대} & \makecell{다세대} & \makecell{준주거} & \makecell{세로\newline{}(가)} & \makecell{정방형\newline{}평지} & \makecell{1,050,000} & \makecell{1,160,000} \\
\hline
\makecell{C} & \makecell{W시 N구\newline{}M동\newline{}101번지} & \makecell{450} & \makecell{대} & \makecell{주상용} & \makecell{준주거} & \makecell{소로\newline{}한면} & \makecell{사다리\newline{}평지} & \makecell{1,100,000} & \makecell{1,210,000} \\
\hline
\end{tabular}\par}

5. 시점수정자료(W시 N구 주거지역) \par
{\small\centering
\begin{tabular}{|c|c|c|}
\hline
\makecell{구분} & \makecell{지가변동률(\%)} & \makecell{비고} \\
\hline
\makecell{2016.01.01～2016.12.31} & \makecell{3.257} & \makecell{2016년 12월 누계} \\
\hline
\makecell{2017.01.01～2017.05.31} & \makecell{1.426} & \makecell{2017년 5월 누계} \\
\hline
\makecell{2017.05.01～2017.05.31} & \makecell{0.431} & \makecell{2017년 5월 변동률} \\
\hline
\end{tabular}\par}

※ 2017년 6월 이후의 지가변동률은 현재 미고시인 상태로 직전월인 2017년 5월 지가변동률을 연장적용 하기로 함 \par
6. 개별요인 품등 비교자료 \par
1) 형상 \par
{\small\centering
\begin{tabular}{|c|c|c|c|c|c|c|}
\hline
\makecell{구분} & \makecell{정방형} & \makecell{가장형} & \makecell{세장형} & \makecell{사다리} & \makecell{부정형} & \makecell{자루형} \\
\hline
\makecell{정방형} & \makecell{1.00} & \makecell{1.02} & \makecell{1.00} & \makecell{0.99} & \makecell{0.94} & \makecell{0.89} \\
\hline
\makecell{가장형} & \makecell{0.98} & \makecell{1.00} & \makecell{0.98} & \makecell{0.97} & \makecell{0.92} & \makecell{0.87} \\
\hline
\makecell{세장형} & \makecell{1.00} & \makecell{1.02} & \makecell{1.00} & \makecell{0.99} & \makecell{0.94} & \makecell{0.89} \\
\hline
\makecell{사다리} & \makecell{1.01} & \makecell{1.03} & \makecell{1.01} & \makecell{1.00} & \makecell{0.95} & \makecell{0.90} \\
\hline
\makecell{부정형} & \makecell{1.06} & \makecell{1.09} & \makecell{1.06} & \makecell{1.05} & \makecell{1.00} & \makecell{0.94} \\
\hline
\makecell{자루형} & \makecell{1.12} & \makecell{1.15} & \makecell{1.12} & \makecell{1.11} & \makecell{1.06} & \makecell{1.00} \\
\hline
\end{tabular}\par}

※ 부정형 : 삼각형 포함 \par
※ 자루형 : 역삼각형 포함 \par
2) 도로접면 \par
{\small\centering
\begin{tabular}{|c|c|c|c|c|c|}
\hline
\makecell{구분} & \makecell{중로한면} & \makecell{소로한면} & \makecell{세로(가)} & \makecell{세로(불)} & \makecell{맹지} \\
\hline
\makecell{중로한면} & \makecell{1.00} & \makecell{0.92} & \makecell{0.82} & \makecell{0.78} & \makecell{0.70} \\
\hline
\makecell{소로한면} & \makecell{1.09} & \makecell{1.00} & \makecell{0.89} & \makecell{0.85} & \makecell{0.79} \\
\hline
\makecell{세로(가)} & \makecell{1.22} & \makecell{1.12} & \makecell{1.00} & \makecell{0.95} & \makecell{0.88} \\
\hline
\makecell{세로(불)} & \makecell{1.28} & \makecell{1.18} & \makecell{1.05} & \makecell{1.00} & \makecell{0.93} \\
\hline
\makecell{맹지} & \makecell{1.43} & \makecell{1.27} & \makecell{1.13} & \makecell{1.07} & \makecell{1.00} \\
\hline
\end{tabular}\par}

7. 그 밖의 요인 보정치 산정을 위한 자료 \par
1) 보상 평가사례 \par
{\small\centering\footnotesize\setlength{\tabcolsep}{3pt}
\begin{tabular}{|c|c|c|c|c|c|c|c|c|c|}
\hline
\makecell{기호} & \makecell{소재지} & \makecell{면적\newline{}(m\textsuperscript{2})} & \makecell{지목} & \makecell{이용상황} & \makecell{용도\newline{}지역} & \makecell{도로\newline{}교통} & \makecell{형상\newline{}지세} & \makecell{단가\newline{}(원/m\textsuperscript{2})} & \makecell{가격시점} \\
\hline
\makecell{ㄱ} & \makecell{W시 N구\newline{}M동\newline{}200번지} & \makecell{300} & \makecell{대} & \makecell{주거\newline{}나지} & \makecell{일반\newline{}상업} & \makecell{소로한면} & \makecell{부정형\newline{}평지} & \makecell{2,500,000} & \makecell{2017.01.01} \\
\hline
\makecell{ㄴ} & \makecell{W시 N구\newline{}M동\newline{}250번지} & \makecell{350} & \makecell{대} & \makecell{다세대} & \makecell{준주거} & \makecell{세로\newline{}(가)} & \makecell{세장형\newline{}평지} & \makecell{1,500,000} & \makecell{2017.01.01} \\
\hline
\makecell{ㄷ} & \makecell{W시 N구\newline{}M동\newline{}300번지} & \makecell{380} & \makecell{대} & \makecell{주상용} & \makecell{준주거} & \makecell{소로한면} & \makecell{사다리\newline{}평지} & \makecell{1,800,000} & \makecell{2017.01.01} \\
\hline
\end{tabular}\par}

2) 거래사례(토지만의 정상 거래사례임) \par
{\small\centering\footnotesize\setlength{\tabcolsep}{3pt}
\begin{tabular}{|c|c|c|c|c|c|c|c|c|c|}
\hline
\makecell{기호} & \makecell{소재지} & \makecell{면적\newline{}(m\textsuperscript{2})} & \makecell{지목} & \makecell{이용상황} & \makecell{용도\newline{}지역} & \makecell{도로\newline{}교통} & \makecell{형상\newline{}지세} & \makecell{단가\newline{}(원/m\textsuperscript{2})} & \makecell{거래시점} \\
\hline
\makecell{ㄹ} & \makecell{W시 N구\newline{}M동\newline{}400번지} & \makecell{400} & \makecell{대} & \makecell{주상\newline{}나지} & \makecell{2종\newline{}일주} & \makecell{중로한면} & \makecell{부정형\newline{}평지} & \makecell{1,600,000} & \makecell{2017.01.01} \\
\hline
\makecell{ㅁ} & \makecell{W시 N구\newline{}M동\newline{}420번지} & \makecell{380} & \makecell{대} & \makecell{상업\newline{}나지} & \makecell{준주거} & \makecell{중로한면} & \makecell{세장형\newline{}평지} & \makecell{2,200,000} & \makecell{2017.01.01} \\
\hline
\makecell{ㅂ} & \makecell{W시 N구\newline{}M동\newline{}500번지} & \makecell{350} & \makecell{대} & \makecell{주거\newline{}나지} & \makecell{2종\newline{}일주} & \makecell{세로\newline{}(가)} & \makecell{사다리\newline{}평지} & \makecell{1,000,000} & \makecell{2017.01.01} \\
\hline
\end{tabular}\par}

8. 현장조사내용 \par
1) 대상토지 주변은 도심지내 기존 주택지를 중심으로 형성된 소규모 점포주택과 단독주택 및 다세대주택 등이 혼재하는 지역으로 조사되었음 \par
2) 비교표준지, 보상 평가사례 및 거래사례는 인근지역에 소재하며, 당해 도시계획시설(도로)사업에 따른 개발이익이 포함되어 있지 않은 것으로 조사되었음 \par
9. 대상 토지 주변 지적현황(축척 없음) \par
\qfigure[0.8\linewidth]{figures/01-1.png}
10. <공통 자료> 외에는 아래 각 물음의 자료를 활용하여 기술하되, 물음 1), 물음 2), 물음 3)은 각각 독립적인 사실관계임을 전제함 \par
<물음 1) 관련 자료> \par
1. 사업시행자인 N구청장은 대상 토지가 미지급용지임을 보상평가의뢰서에 명기하였음 \par
2. 본건 토지는 종전 공익사업의 시행으로 M동 100번지에서 분필된 토지로 현재 도로로 이용중인 사다리형 토지이나, 종전 편입당시에는 부정형, 맹지인 토지였음 \par
3. 본건 토지는 종전 공익사업 시행이전에는 제2종일반주거지역 이었으나, 종전의 공익사업으로 인하여 준주거지역으로 변경되었음 \par
4. 본건 토지의 지목은 종전의 공익사업에 편입되기 전에도 “전” 이었으나, 당시 주변 토지의 표준적 이용상황 등을 고려할 때 종전 편입당시의 이용상황은 주거나지였음을 사업시행자로부터 확인하였음 \par
5. 개별요인 비교시 도로접면과 형상을 제외한 토지 특성은 모두 동일한 것으로 봄 \par
6. 비교표준지, 보상 평가사례 및 거래사례에는 종전 공익사업에 따른 가치변동이 포함되어 있지 아니함 \par
<물음 2) 관련 자료> \par
1. 사업시행자인 N구청장은 대상 토지가 사실상 사도임을 보상평가의뢰서에 명기하였음 \par
2. 인접한 M동 100-2번지 소유권자는 대상 토지와 동일한 홍길동이고, 1992년 8월부터 홍길동은 M동 100-2번지의 건축허가를 위하여 대상 토지를 도로로 개설한 것으로 확인이 되었는바, 대상 토지는 M동 100-2번지의 효용증진에 기여하고 있음 \par
3. 대상 토지 평가시 기준이 되는 인근토지 및 인근토지의 토지 특성은 주어진 자료를 활용하여 판단할 것 \par
4. 개별요인 비교시 도로접면과 형상을 제외한 토지 특성은 모두 동일한 것으로 봄 \par
<물음 3) 관련 자료> \par
1. 사업시행자인 N구청장은 대상 토지가 예정공도임을 보상평가의뢰서에 명기하였음 \par
2. 인접한 M동 100-2번지 소유권자는 대상 토지와 동일한 홍길동이며, 당해 도시계획시설사업 시행절차 등이 없는 상태에서 M동 100-2번지의 건축허가를 위하여 2012년 8월부터 도로로 개설한 후, 사실상 불특정 다수인의 통행에 이용 중임을 사업시행자로부터 확인하였음 \par
3. 인근지역의 표준적인 이용상황은 주어진 자료를 활용하여 판단할 것 \par
4. 개별요인 비교시 도로접면과 형상을 제외한 토지 특성은 모두 동일한 것으로 보며, 당해 도로의 개설로 인한 개발이익은 30\%임 \par
5. 주변토지 이용상황 \par
{\small\centering
\begin{tabular}{|c|c|c|c|c|c|}
\hline
\makecell{구분} & \makecell{소재지} & \makecell{지목} & \makecell{이용상황} & \makecell{용도지역} & \makecell{비고} \\
\hline
\makecell{가} & \makecell{W시 N구 M동\newline{}104-1, 104-14번지} & \makecell{대} & \makecell{단독주택} & \makecell{2종일주} & \makecell{-} \\
\hline
\makecell{나} & \makecell{W시 N구 M동\newline{}104-2, 104-3번지} & \makecell{도} & \makecell{도로} & \makecell{2종일주} & \makecell{-} \\
\hline
\makecell{다} & \makecell{W시 N구 M동\newline{}109, 109-1번지} & \makecell{대} & \makecell{단독주택} & \makecell{준주거지역} & \makecell{-} \\
\hline
\makecell{라} & \makecell{W시 N구 M동\newline{}110, 110-1번지} & \makecell{대} & \makecell{단독주택} & \makecell{준주거지역} & \makecell{-} \\
\hline
\makecell{마} & \makecell{W시 N구 M동\newline{}102번지} & \makecell{대} & \makecell{단독주택} & \makecell{준주거지역} & \makecell{-} \\
\hline
\makecell{바} & \makecell{W시 N구 M동\newline{}100-2번지} & \makecell{대} & \makecell{다세대} & \makecell{준주거지역} & \makecell{-} \\
\hline
\makecell{사} & \makecell{W시 N구 M동\newline{}100, 100-8번지} & \makecell{대} & \makecell{다세대} & \makecell{준주거지역} & \makecell{-} \\
\hline
\makecell{아} & \makecell{W시 N구 M동\newline{}100-1번지} & \makecell{전} & \makecell{도로} & \makecell{준주거지역} & \makecell{-} \\
\hline
\makecell{자} & \makecell{W시 N구 M동\newline{}100-6, 100-9번지} & \makecell{전} & \makecell{주상나지} & \makecell{준주거지역} & \makecell{-} \\
\hline
\makecell{차} & \makecell{W시 N구 M동\newline{}100-7번지} & \makecell{대} & \makecell{주상용} & \makecell{준주거지역} & \makecell{-} \\
\hline
\makecell{카} & \makecell{W시 N구 M동\newline{}99번지} & \makecell{대} & \makecell{주상용} & \makecell{준주거지역} & \makecell{-} \\
\hline
\end{tabular}\par}

\end{longquestion}
